\section{An infinite-order coarse K-homology class}\label{sec:persistence}

The small-loop lemma converts winding on the base lattice into an integral
cocycle on each sufficiently late Rips component. We use this cocycle to
test the same boundary family at every scale.

All links between components of $B$ have length greater than one, and each
$C_j$ has length $6j\ge12$. Thus $P_1(B)$ is exactly the topological disjoint
union of the circles $C_j$. By \eqref{eq:additivity}, their positively
oriented analytic classes determine
\[
 \beta_1=([C_j])_{j\ge2}\in K_1^{\lf}(P_1(B)).
\]
Let $\beta\in KX_1(B)$ be its direct-limit image. For each integer $r\ge1$,
let $i_r:P_1(B)\to P_r(X)$ be the map induced by the boundary inclusion,
and put
\[
 \alpha_r=(i_r)_*\beta_1,\qquad \alpha=i_*\beta\in KX_1(X).
\]
The classes $\alpha_r$ represent one and the same element $\alpha$.

\begin{proposition}\label{prop:persistence}
For every integer $r\ge1$ there is a homomorphism
\[
 \lambda_r:K_1^{\lf}(P_r(X))\longrightarrow\Z
 \quad\text{such that}\quad \lambda_r(\alpha_r)=1.
\]
Consequently, every nonzero integer multiple of $\alpha$ is nonzero in
$KX_1(X)$.
\end{proposition}
\begin{proof}
Choose $j$ sufficiently large that $j\ge J(3r)$ and that both auxiliary
links adjacent to $X_j$ have length greater than $r$. Then $Y=P_r(X_j)$ is
a finite clopen component of $P_r(X)$.

For an oriented edge $(x,y)$ of $Y$, choose an edge path in $X_j$ from $x$
to $y$ of length at most $r$, and let $b_r(x,y)$ be its charge. Two choices
differ by a closed path of length at most $2r$, so
Lemma~\ref{lem:smallloops} makes this integer independent of the choice.
Reversing a path gives $b_r(y,x)=-b_r(x,y)$. Around any Rips triangle, the
three chosen paths form a closed path of length at most $3r$, again with
zero charge. Therefore $b_r$ is an integral simplicial $1$-cocycle on $Y$.

On each unit boundary edge of $C_j$ we may choose that edge itself. Hence
the sum of $b_r$ around the positively oriented $C_j$ is one. By the
construction in Section~\ref{sec:conventions}, there is a map
$f_r:Y\to S^1$ whose unitary class pairs with the included boundary as
\[
 \langle[f_r],(C_j\hookrightarrow Y)_*[C_j]\rangle=1.
\]
Define $\lambda_r$ by restriction of locally finite K-homology to the
clopen component $Y$, followed by this index pairing. Restriction of
$\alpha_r$ is exactly the included $j$th boundary class. Indeed,
$i_r^{-1}(Y)=C_j$, and pullback along $i_r$ carries extension by zero from
$C(Y)$ to extension by zero from $C(C_j)$, after restriction to $C_j$.
Applying the contravariant first variable of $\KK$ gives the claimed
coordinate projection in \eqref{eq:additivity}, regardless of connections
among the other Rips components. This proves the displayed claim.
Equivalently, one may extend $f_r$ by $1$ outside $Y$ to obtain a
unitary in the unitization of $C_0(P_r(X))$; it is a compactly supported
K-theory test for the locally finite class.

For any integer $n\ne0$, we have $\lambda_r(n\alpha_r)=n\ne0$ at every
integer scale $r\ge1$. If $n\alpha$ vanished in the direct limit, its
representative $n\alpha_1$ would vanish at some later scale, and then at
some integer scale. This contradicts the pairing. Integer scales are
cofinal, so the conclusion holds in $KX_1(X)$.
\end{proof}

The component and the compactly supported test may change with $r$. The
argument requires one test at each scale, and does not assert a single
test persisting at all scales. We have proved that $\alpha$ has infinite
order integrally. Since the kernel of $A\to A\otimes_{\Z}\Q$ is the
torsion subgroup for any abelian group $A$, it follows that
\begin{equation}\label{eq:rationalnonzero}
 \alpha\otimes1\ne0\quad\text{in }KX_1(X)\otimes_{\Z}\Q.
\end{equation}
No product in \eqref{eq:additivity} has been rationalized separately.
